% Sommario (abstract) dell'allegato: solo il testo, da racchiudere in \begin{abstract}...\end{abstract} o in un paragrafo «Sommario».
Si descrive un sistema di percezione che dichiara ciò che non ha osservato. Per ogni oggetto tiene distinte tre risposte: lo vede; non lo vede, ma indica una regione dove può essere; non sa. Il non osservato non è mai dichiarato «libero»; l'avviso di possibile presenza si spegne solo dopo aver guardato tutta la zona; la perdita di affidabilità di un sensore è dichiarata. Il caso centrale è il volo ad alta velocità di aerei e droni, da $20$ a $210\unita{m/s}$, con collisioni, detriti ed eventi improvvisi fortemente impulsivi. Tutto è simulazione dalle schede tecniche, calcolo, o traiettorie pubbliche con un sensore simulato: nessuna prova usa un sensore fisico o hardware dedicato, il classificatore è supposto esistente e le etichette sono assegnate a mano, e il sistema fa solo osservazione, monitoraggio e avviso. Il prototipo è stato costruito e provato in un giorno.

Nei sei scenari di volo simulati, dall'avaria di un aereo alla collisione di due droni veloci, il sistema non ha mai dichiarato «libero» dove c'era l'oggetto o un frammento ($0/200$ $[0;\,0{,}019]$ nel controllo, $0/60$ $[0;\,0{,}060]$ nelle altre celle; intervalli di Wilson al $95\,\%$ per passaggio). Dopo la collisione veloce il punto vero e i frammenti sono rimasti nella regione dichiarata in $60/60$ passaggi $[0{,}940;\,1]$, grazie a un insieme di presenza con contenimento dimostrato per velocità fino a $210\unita{m/s}$, se l'errore iniziale è entro $r_0$. La collisione veloce, invece, non è ancora dichiarata come evento: una regola di contatto in prova la riconosce entro $2\unita{s}$ in 57 passaggi su 60 $[0{,}863;\,0{,}983]$ con un radar a scansione elettronica, ma apre ancora falsi eventi in un quasi incrocio (2 su 60) e non è adottata. A bordo di un drone in moto, durante una collisione a $30$--$100\unita{m}$ con fiammata, fumo e detriti, la regione del drone seguito non si è persa in $60/60$ passaggi per geometria con un radar a $360^\circ$; con il radar accecato fino a $2\unita{s}$ il «libero» resta assente, ma l'avviso dei detriti al portatore non è affidabile.

La perdita di confidenza dopo l'ultima misura è data in forma chiusa, $C(\tau)=1-\exp(-r^2/2\sigma^2(\tau))$ con $\sigma^2(\tau)$ polinomio cubico in $\tau$, dimostrata nel modello a velocità quasi costante ed estesa a una miscela di regimi di moto. Dopo una taratura, su traiettorie di volo pubbliche con un sensore simulato, la regione al $90\,\%$ ha contenuto l'aereo da 85 a 89 volte su 100 fra $1$ e $30\unita{s}$; i casi vengono per nove decimi da quattro voli di una sola rotta, e in virata la copertura scende a $0{,}73$ $[0{,}67;\,0{,}78]$, un limite dichiarato. Sensori allo stato dell'arte e calcolo dedicato sono stimati per analogia: il calcolo dedicato porterebbe l'elaborazione a $0{,}08$--$1{,}2\unita{ms}$ (ipotesi), ma la leva vera è la rivisita del sensore. Ogni limite trovato ha la sua miglioria e la prova che la deciderebbe.

Gli altri usi condividono lo stesso nucleo. A corto raggio, un prototipo rileva elementi mobili e ostacoli con una camera di profondità a stereo attivo, frontale e ferma, e continua a dare un'indicazione dichiarata quando un ostacolo nasconde il soggetto. Il sistema tiene distinte tre informazioni: l'osservato; il non osservato ma previsto, con una regione e un livello; il non noto. Lo stato «libero» è scritto solo per il suolo osservato nel fotogramma corrente; l'avviso di possibile presenza non termina senza osservazione diretta. Nelle ipotesi dichiarate sono dimostrati che ciò che non è osservato non è mai «libero», il contenimento del soggetto nascosto nell'insieme di possibile presenza (per una lettura esatta del suolo, senza il margine che il sistema usa) e la condizione sulla fascia di guardia. In simulazione, con modello ideale del sensore e un soggetto alla volta, la regione di previsione è emessa a ogni ciclo dell'occlusione una volta riconosciuto l'ingresso nell'ombra (0 interruzioni dopo l'ingresso su $23\,402$ passaggi nascosti; in 928 di essi la traccia si chiude prima dell'ombra e non c'è regione, e con una regola aggiuntiva, attiva solo dopo una prova d'ingresso, 24 restano senza avviso); con un ostacolo fermo largo $2\unita{m}$, il bordo dell'ombra noto e il soggetto sul proprio percorso, la regione contiene il soggetto in almeno 9 casi su 10 fino a $2\unita{s}$ e in almeno 85 su 100 fino a $3\unita{s}$; con un ostacolo largo $1\unita{m}$ la soglia di 9 su 10 non è raggiunta, e su profondità simulata dalle schede tecniche questa regola non è stata eseguita. Dietro un ostacolo fermo l'avviso resta acceso, e nella realizzazione provata resta acceso anche dopo la rimozione dell'ostacolo. Il classificatore degli oggetti è un componente richiesto dal sistema: in questo lavoro non è realizzato né valutato, ed è simulato con etichette assegnate a mano. Nessuna prova è stata eseguita su un sensore fisico: i risultati a corto raggio sono ottenuti a camera ferma, al chiuso, fino a $3\unita{m}$, e non sono stati verificati fuori da queste condizioni.

Su un sito aerospaziale sorvegliato, con droni e aerei in rullaggio, i sensori a lunga portata sono simulati dalle schede tecniche (radar per droni, radar del movimento a terra, LiDAR, camere, ADS-B e identificazione remota). Con $3\unita{m}$ di portata nessuno di questi scenari è servito: la distanza d'avviso richiesta va da $24\unita{m}$ a qualche chilometro. Il «libero» si dichiara per classe, con la probabilità di rilevamento cumulata e le passate trattate come correlate; l'avviso usa l'inviluppo dello scenario e si spegne solo guardando tutta la zona in cui il mezzo può essere. In simulazione dalle schede, un drone a 15--$21\unita{m/s}$ nascosto da un edificio è annunciato in tempo all'operatore in 200 passaggi su 200, e un aereo in rullaggio nascosto da un hangar in 200 su 200; una regola di spegnimento più semplice perde l'avviso in 564 passaggi su 2400. Per gli aerei in volo l'avviso non è tempestivo, e il sistema lo dichiara.
Le simulazioni mostrano anche dove il sistema cede. Il «libero» dipende dalla mappa degli ostacoli: con la mappa del sito dichiarata completa (un'ipotesi del simulatore) nessuno di 171 passaggi ha un «libero» sul drone, ma con una mappa dichiarata e sbagliata il «libero» dietro un edificio che non vi compare è falso; per questo, per impostazione, la mappa non è dichiarata completa. Quando il drone rallenta o si ferma l'avviso lo perde per qualche ciclo in 12 passaggi su 200. Dopo una collisione fra droni lenti l'evento è dedotto entro $2\unita{s}$ in 48 passaggi su 60 con l'integratore stabile dei frammenti (51 con il primo, instabile), con 3 eventi aperti prima dell'urto, e la zona dei detriti è di alcuni chilometri e non si riduce guardando. Le probabilità di rilevamento e le latenze dei sensori sono ipotesi.
